\subsection{平方根完成化坐标与零温饱和缺口}\label{subsec:conclusion-squareroot-saturation-completed-coordinate}
沿用命题 \ref{prop:real-input-40-pressure-freezing}、定理 \ref{thm:xi-output-potential-zero-temp-square-root-critical-law}、命题 \ref{prop:weighted-completion-Q} 与定理 \ref{thm:conclusion-selfdual-sync-kernel-finite-mirror-odd-cumulant} 的记号。前述链条一方面已经把真实输入 \(40\) 状态核在零温端点的 Perron Puiseux jet 固定为平方根分支，另一方面又把自对偶同步核的有限长度配分多项式锁定在回文不变环中。于是平方根完成化并非只存在于无穷长的 Perron 分支，而同时控制饱和缺口与每一个有限长度配分函数的严格降阶。

\begin{theorem}[\texorpdfstring{$40$}{40} 状态真实输入核的饱和缺口精确指数律]\label{thm:conclusion-realinput40-saturation-gap-halfpower}
沿用命题 \ref{prop:real-input-40-pressure-freezing} 的记号。记
\[
u=e^\theta,
\qquad
P(\theta)=\log\lambda(u),
\qquad
\alpha(u):=\frac{u\lambda'(u)}{\lambda(u)}=P'(\theta),
\]
并令 \(c,d\) 为命题 \ref{prop:real-input-40-pressure-freezing} 中的 Puiseux 系数。则有精确饱和缺口律
\[
\frac12-\alpha(u)=\frac{d}{2c}\,u^{-1/2}+O(u^{-1})
\qquad (u\to+\infty),
\]
等价地，
\[
\frac12-P'(\theta)=\frac{d}{2c}e^{-\theta/2}+O(e^{-\theta})
\qquad (\theta\to+\infty).
\]
其中
\[
c\approx 1.89635299351377,
\qquad
d\approx 0.807943317323508,
\]
故前导常数为
\[
\frac{d}{2c}\approx 0.2130255601.
\]
\end{theorem}
\begin{proof}
由命题 \ref{prop:real-input-40-pressure-freezing}，
\[
P(\theta)
=
\frac{\theta}{2}+\log c+\frac{d}{c}e^{-\theta/2}
+\Bigl(\frac{e}{c}-\frac{d^2}{2c^2}\Bigr)e^{-\theta}
+O(e^{-3\theta/2})
\qquad (\theta\to+\infty).
\]
逐项求导即得
\[
P'(\theta)
=
\frac12-\frac{d}{2c}e^{-\theta/2}
-\Bigl(\frac{e}{c}-\frac{d^2}{2c^2}\Bigr)e^{-\theta}
+O(e^{-3\theta/2}),
\]
从而
\[
\frac12-P'(\theta)=\frac{d}{2c}e^{-\theta/2}+O(e^{-\theta}).
\]
再由 \(u=e^\theta\) 与 \(\alpha(u)=P'(\theta)\) 代回，即得
\[
\frac12-\alpha(u)=\frac{d}{2c}u^{-1/2}+O(u^{-1}).
\]
这也正是定理 \ref{thm:xi-output-potential-zero-temp-square-root-critical-law} 的第一条结论在 \(\alpha(u)\) 记号下的重述。证毕。
\end{proof}

\begin{proposition}[有限长度配分函数的唯一降阶坐标]\label{prop:conclusion-sync-kernel-finite-length-single-coordinate}
\leanverified{paper\_conclusion\_sync\_kernel\_finite\_length\_single\_coordinate}
在定理 \ref{thm:conclusion-selfdual-sync-kernel-finite-mirror-odd-cumulant} 的设定下，对每个 \(n\ge 1\)，存在唯一多项式
\[
Q_n(T)\in\mathbb Z[T]
\]
使得
\[
Z_n(u)=u^{n/2}Q_n\!\left(u^{1/2}+u^{-1/2}\right).
\]
等价地，若令 \(u=e^\theta\)，则存在唯一 \(Q_n\) 满足
\[
Z_n(e^\theta)=e^{n\theta/2}Q_n\!\left(2\cosh\frac{\theta}{2}\right).
\]
在单位圆扭曲 \(u=e^{i\phi}\) 上，上式化为
\[
Z_n(e^{i\phi})=e^{in\phi/2}Q_n\!\left(2\cos\frac{\phi}{2}\right).
\]
因此，自对偶同步核在每一个有限长度上的 canonical family 实际上都只通过单一实坐标
\[
T=u^{1/2}+u^{-1/2}
\]
传递；所有有限长度的输出统计在代数上已完成一次严格降维。
\end{proposition}
\begin{proof}
由定理 \ref{thm:conclusion-selfdual-sync-kernel-finite-mirror-odd-cumulant}，
\[
Z_n(u)=u^n Z_n(1/u).
\]
令 \(r:=u^{1/2}\)，并置
\[
\widetilde Z_n(r):=r^{-n}Z_n(r^2).
\]
则
\[
\widetilde Z_n(r)=\widetilde Z_n(r^{-1}),
\]
故 \(\widetilde Z_n\) 落在不变环
\[
\mathbb Z[r,r^{-1}]^{\,r\mapsto r^{-1}}.
\]
而
\[
\mathbb Z[r,r^{-1}]^{\,r\mapsto r^{-1}}=\mathbb Z[r+r^{-1}],
\]
故存在唯一 \(Q_n\in\mathbb Z[T]\) 使
\[
\widetilde Z_n(r)=Q_n(r+r^{-1}).
\]
把 \(r=u^{1/2}\) 代回，即得
\[
Z_n(u)=u^{n/2}Q_n\!\left(u^{1/2}+u^{-1/2}\right).
\]
取 \(u=e^\theta\) 时，\(u^{1/2}+u^{-1/2}=2\cosh(\theta/2)\)；取 \(u=e^{i\phi}\) 时，\(u^{1/2}+u^{-1/2}=2\cos(\phi/2)\)，从而得到两条等价表述。唯一性则来自嵌入
\[
\mathbb Z[T]\hookrightarrow \mathbb Z[r,r^{-1}],
\qquad
T\mapsto r+r^{-1}
\]
的单射性。证毕。
\end{proof}

\endinput
